\chapter{Eine Bitte}
\quellseite{11}

Im Anschluß an diese Nominationen erlaubt sich der Unterzeichnete, einem Wunsche Ausdruck zu geben. Bei der großen Zahl unserer Anstalten und darum unserer Professoren tritt die ernste Pflicht der Wahl an die Glieder des Wahlcollegiums immer häufiger heran. Diese sollen gewissenhaft den Namen des nach ihrer Ueberzeugung für das betreffende Amt tüchtigsten Candidaten aus der Wahlurne hervorgehen lassen. Um dazu im Stande zu sein, müssen sie über die in Frage kommenden Personen hinreichende Information besitzen, damit sie sich ein Urtheil bilden können. Nun hat zwar die Synode bei der Zusammenstellung des Wahlcollegiums auf das Synodalgebiet in diesem Lande Rücksicht genommen und die Glieder ihrer Wahlbehörde so herausgesucht, daß jeder Theil möglichst vertreten ist. Bei der Menge der Pastoren und Lehrer, aus deren Zahl die Candidaten genommen werden, und bei der großen geographischen Ausbreitung unserer Synode tritt jedoch nicht selten der Fall ein, daß über Befähigung etc. des einen oder mehrerer Candidaten kein gewisses Zeugniß vorliegt, zumal die meisten Nominationenanzeigen nur die nackte Thatsache der Nomination bringen. Die aber, welchen die Vollziehung der Wahlen obliegt, besitzen aus dem angegebenen Grunde nicht die erwünschte Personalkenntniß aller Candidaten. Das ist unmöglich und wird je länger je mehr unmöglich werden. Wer daher jetzt oder bei späteren Wahlen über Gaben, Fertigkeit, und zwar besonders in den zu lehrenden Fächern, Lehrtüchtigkeit etc. eines oder mehrerer der Nominierten sachlichen Aufschluß nach der einen oder der anderen Seite hin zu geben vermag, wird der Kirche einen nicht zu unterschätzenden Dienst leisten, wenn er mündlich, lieber schriftlich, irgend einem Gliede des Wahlcollegiums Mittheilung macht. Solche Zeugnisse sind selbstverständlich nur für die Wahlbehörde bestimmt und werden als solche sub sigillo gehalten werden.

\begin{unterschrift}
L. Hölter, d. Z. Vorsitzender
\end{unterschrift}
